% Compile with pdflatex (twice). This file needs no external bibliography.
\documentclass[10pt,a4paper]{article}
\usepackage[margin=18mm]{geometry}
\usepackage[T1]{fontenc}
\usepackage{lmodern}
\usepackage{array,booktabs,tabularx}
\usepackage[hidelinks]{hyperref}
\setlength{\parindent}{0pt}
\setlength{\parskip}{4pt}
\newcommand{\topic}[1]{\par\medskip\textbf{#1}\par\smallskip}
\newcommand{\todo}[1]{\textit{[To fill: #1]}}
\begin{document}
\begin{center}
  {\Large\bfseries Runtime Repair of Game Automation Workflows\\with LLM Agents}\par
  \vspace{3pt}
  NLP Final Project Proposal
\end{center}
\textbf{Team:} \todo{names, student IDs, and course codes}.\\
\textbf{Track:} A (Agents \& Harness), custom study related to A4 (bounded recovery).

\topic{1. Question and motivation}
At a fixed model-request budget, how do game-state knowledge, workflow-repair
guidance, and reference repositories affect an LLM agent's ability to repair a
failing game-automation workflow and complete its task? We study these resources
in LALC, an existing system that lets a coding agent inspect screenshots and
execution feedback, edit workflow code, and submit revisions for executor-controlled
resumption. The contribution is a controlled evaluation of runtime repair, rather
than workflow generation, model training, or a claim of a new agent algorithm.

\topic{2. Closest prior work}
Voyager uses execution feedback to improve executable game skills~\cite{voyager}.
SWE-agent studies interfaces for interactive code repair~\cite{sweagent}.
AndroidWorld emphasizes task initialization and independent success
checks~\cite{androidworld}. We instead evaluate repair of an existing Python
workflow during execution in a real desktop game, and isolate the contribution
of three reference resources under a shared request allowance.

\topic{3. Environment, cases, and held-out evaluation}
In \emph{Limbus Company}, the fixed task is to start at the main menu, enter the
first floor of a dungeon (the game's ``Mirror Dungeon''), leave the dungeon, and
return to the main menu. No full dungeon completion is required. A verified clean
workflow serves as the starting artifact. Ten candidate single-fault cases are
listed below; the supplied version IDs are abbreviated for readability.

\begingroup
\small
\renewcommand{\arraystretch}{1.12}
\begin{tabularx}{\linewidth}{@{}l l X@{}}
\toprule
Case & Version prefix & Injected fault or changed condition \\
\midrule
% 22fae48e-9f41-48aa-8102-e8c3b2f0b1b1
C1 & \texttt{22fae48e} & Read \texttt{repetitions["tests"]}: wrong key, \texttt{KeyError}. \\
% 7193ac2f-6d68-48e0-9c67-b270159df771
C2 & \texttt{7193ac2f} & Concatenate a string and numeric \texttt{total}: \texttt{TypeError}. \\
% d999b5a3-8bec-4c12-912a-15452aa0fb68
C3 & \texttt{d999b5a3} & Log undefined \texttt{total\_count}: \texttt{NameError}. \\
% 8b5fa5b5-ddf4-4d66-b35f-b53b9af3aea1
C4 & \texttt{8b5fa5b5} & Index \texttt{SUPPORTED\_TASK\_IDS[1]}: \texttt{IndexError}. \\
% e77c8b6c-45a2-437a-be7d-502ac05f18e4
C5 & \texttt{e77c8b6c} & Call \texttt{.strip()} on a numeric repetition count: \texttt{AttributeError}. \\
% 86197633-024e-44d5-ad00-39f66f1398a9
C6 & \texttt{86197633} & Reject valid counts with \texttt{total >= 0}: faulty validation logic. \\
% cd8e2552-0efe-4b36-b7dc-5dc6e3f788e5
C7 & \texttt{cd8e2552} & Read local \texttt{phase} before assignment: \texttt{UnboundLocalError}. \\
% 6d96202d-36b0-409c-8754-f0c455297bd3 (also the clean workflow)
C8 & \texttt{6d96202d} & Clean workflow with an icon-position change: stale targeting. \\
% 08c3f7b2-0f84-4027-9d37-3a0b4af3e4c4
C9 & \texttt{08c3f7b2} & Icon-appearance change: stale visual recognition. \\
% ea2b202a-b244-496d-ac0e-ca0e04bd2c7a
C10 & \texttt{ea2b202a} & Remove module-level \texttt{\_main\_evidence()}: missing callable. \\
\bottomrule
\end{tabularx}
\endgroup

Before freezing the benchmark, we will verify clean execution and each fault's
reproducible runtime manifestation. Visual cases are identified by workflow
\emph{and} UI/asset condition; C8 shares its code version with the clean control.
Separate pilot faults will be used for setup. The ten evaluation cases, prompts,
resource snapshots, and scoring rubric will then be frozen; no tuning on their
results is allowed. Clean source, injected diffs, and reference fixes are hidden
from the agent, while ordinary failing code and runtime feedback remain visible.

\topic{4. Baseline and 5. Targeted change}
The baseline is the existing LALC repair agent without the three optional
reference resources. All conditions retain identical required API contracts,
tools, permissions, and deterministic recovery rules. We test resource
availability, with five conditions: \textbf{Baseline} (none), \textbf{Full} (all),
\textbf{Full--Game}, \textbf{Full--Repair}, and \textbf{Full--Repositories}.
Game knowledge and repair guidance are read-only Skill documents; repositories
are version-pinned, searchable code and documentation from similar automation
projects. We log actual resource reads. Full versus Baseline is the main contrast;
single removals estimate each resource's contribution within Full, not all
interactions. Baseline performance is unknown and will be measured, not assumed.

\newpage
\topic{6. Protocol, scale, and metrics}
Each condition runs once on each of the ten cases:
\textbf{5 conditions $\times$ 10 cases = 50 repair episodes}.
Each episode resets the main-menu state, faulty artifact, and session, without
inherited patches or history. Model version, settings, task instructions, account
configuration, resolution, and assets are fixed across conditions; run order is
balanced. The operator initializes and observes, without repair hints or manual
task completion.

Each episode receives \textbf{at most 100 model requests in total}, shared across
all repair interventions and failed revisions, without resetting on re-entry.
There is \textbf{no episode wall-clock limit}; actual elapsed time is measured
from faulty-workflow startup through completion, abandonment, or request-budget
exhaustion. Manual setup and assessment time are recorded separately.

\textbf{Primary metric:} independently verified end-to-end task success within
the request allowance (successful episodes / evaluable episodes). A reviewer
checks preserved screenshots or video and execution traces against the frozen
rubric: actual entry to the first floor, actual exit, and return to the main menu.
Normal program termination, bypassing the task, or an agent's success claim alone
does not count. The rubric cannot be edited by the agent.

\textbf{Secondary metrics:} requests, input/output tokens, API cost, elapsed time,
and total API cost divided by successes. Report abandonment, budget exhaustion,
and false completion claims (incorrect / all success claims). Show case-level
outcomes and include failed attempts in costs. Document and rerun infrastructure
interruptions under a predeclared rule. Archive configurations, version IDs,
resource reads, patches, and traces.

Prioritizing ten cases over repeated runs follows TA feedback. With one episode
per case and condition, results are exploratory paired comparisons; model
variance and general game competence cannot be established.

\topic{7. Resource budget}
Use one Windows machine, an authorized game account, and an API model; no
training-GPU hours are needed. The main study allows \textbf{5,000 model requests},
excluding separately tracked pilots; token usage determines actual spend.
\textbf{Before launch:} \todo{model/provider, pilot allowance, and API spending cap
based on pilot costs}. Pin repository commits and Skills. Existing LALC
engineering predates this study; benchmark construction, experiments, analysis,
and reproducibility materials constitute the course contribution.

\topic{8. Four-week plan}
\textbf{W1:} verify the clean route and faults, pilot on separate cases, finalize
model and cost cap, and freeze the benchmark/resources/rubric.
\textbf{W2:} run the five-condition matrix with resets and preserved traces.
\textbf{W3:} independently score results, compare paired case outcomes and costs,
and analyze failures and resource use.
\textbf{W4:} audit headline results against logs, reproduce selected cases,
prepare a live or recorded repair demo, and complete the report and poster.

\topic{9. Fallback experiment}
If resources do not improve success, retain the negative result and analyze
whether references were unused, redundant, or unhelpful for particular faults.
If the pilot predicts excessive cost, reduce the predeclared design to Baseline
versus Full on all ten cases (20 episodes), preserving case coverage and the
100-request allowance. If visual perturbations cannot be made reproducible,
freeze a code-fault-only study and explicitly narrow its claims before evaluation.

\topic{10. Risks and teaching-team support}
Simple plan-level faults may produce ceiling effects and reveal little about
game knowledge. Two visual cases cannot establish general update robustness.
Audit reference overlap and exclude benchmark answers. Control game updates,
account state, and accidental resource spending. We seek feedback on case
diversity and analysis scope. The TA has confirmed the topic, scoped task, and
initial scale, prioritizing additional cases over repetitions.

\begingroup
\small
\setlength{\parskip}{0pt}
\begin{thebibliography}{3}
\setlength{\itemsep}{0pt}
\bibitem{voyager} G. Wang et al. (2023).
\emph{Voyager: An Open-Ended Embodied Agent with Large Language Models}.
\href{https://arxiv.org/abs/2305.16291}{arXiv:2305.16291}.
\bibitem{sweagent} J. Yang et al. (2024).
\emph{SWE-agent: Agent-Computer Interfaces Enable Automated Software Engineering}.
\href{https://arxiv.org/abs/2405.15793}{arXiv:2405.15793}.
\bibitem{androidworld} C. Rawles et al. (2024).
\emph{AndroidWorld: A Dynamic Benchmarking Environment for Autonomous Agents}.
\href{https://arxiv.org/abs/2405.14573}{arXiv:2405.14573}.
\end{thebibliography}
\endgroup
\end{document}
